%==============================================================
% ABSTRAK DAN ABSTRACT
%==============================================================

\clearpage
\thispagestyle{empty}

% --- ABSTRAK (Bahasa Indonesia) ---
\begin{center}
{\fontsize{14pt}{17pt}\selectfont\bfseries ABSTRAK}
\end{center}

\vspace{6pt}

\noindent
GHIFFARI BRAVIA HISHAM, MUHAMMAD ABYAN PUTRA WIBOWO, ADAM NAUFAL,
MOCHAMAD CHAIRULRIDJAL NURVIKRI, RAFIF MUHAMMAD FARRAS.
Klasifikasi \textit{Stage Retinopathy of Prematurity} Berdasarkan Citra Fundus Bayi
Menggunakan Representasi \textit{Softmap} Spasial dan \textit{Custom} CNN.
Mata kuliah Pengenalan Citra Digital. Dibimbing oleh NAMA DOSEN.

\vspace{12pt}

\noindent
\textit{Retinopathy of Prematurity} (ROP) adalah gangguan perkembangan pembuluh darah retina
pada bayi prematur yang berpotensi menyebabkan kebutaan jika tidak terdeteksi sejak dini.
Penelitian ini membangun sebuah \textit{pipeline} klasifikasi \textit{stage} ROP dari citra
fundus bayi ke lima kelas: Normal, Stage 1, Stage 2, Stage 3, dan bekas luka laser
(\textit{laser scars}). Penelitian membandingkan citra RGB mentah dengan representasi
\textit{softmap} spasial yang dirancang untuk menonjolkan struktur klinis penting. Model
klasifikasi utama adalah arsitektur CNN residual kecil (\textit{TinyResNet}) yang dilatih dari
awal tanpa bobot pra-latih, menerima masukan tiga kanal berupa \textit{softmap} pembuluh darah
berbasis Gabor, \textit{softmap} area tepi berbasis Meijering dan \textit{oriented top-hat}, serta
kanal hijau ternormalisasi CLAHE. Validasi silang 5-\textit{fold} yang
mempertimbangkan kelompok citra serupa (\textit{group-aware cross-validation}) menghasilkan
\textit{macro F1-score} sebesar \textbf{0,7853}. Sebagai evaluasi pendukung, alur segmentasi
pembuluh darah mencapai nilai Dice \textbf{0,3881} pada dataset HVDROPDB-BV, sedangkan kanal
\textit{ridge} mencapai Dice \textbf{0,0993} pada HVDROPDB-RIDGE. Percobaan ablasi yang terkontrol menunjukkan bahwa
representasi \textit{softmap} memberikan peningkatan 0,0987 dibanding masukan RGB mentah
(0,6866 vs 0,7853) pada model, data, dan protokol yang sama. Hasil ini membuktikan bahwa
representasi spasial dari informasi pembuluh darah---bukan hanya nilai piksel mentah---merupakan
kunci utama peningkatan performa klasifikasi \textit{stage} ROP.

\vspace{12pt}

\noindent
Kata kunci: \textit{custom} CNN, citra fundus, \textit{Retinopathy of Prematurity},
representasi \textit{softmap}, \textit{TinyResNet}

\vspace{1.5cm}

% --- ABSTRACT (English) ---
\begin{center}
{\fontsize{14pt}{17pt}\selectfont\bfseries ABSTRACT}
\end{center}

\vspace{6pt}

\noindent
GHIFFARI BRAVIA HISHAM, MUHAMMAD ABYAN PUTRA WIBOWO, ADAM NAUFAL,
MOCHAMAD CHAIRULRIDJAL NURVIKRI, RAFIF MUHAMMAD FARRAS.
Classification of Retinopathy of Prematurity Stages from Infant Fundus Images
Using Spatial Softmap Representation and TinyResNet.
Digital Image Recognition Course. Supervised by NAMA DOSEN.

\vspace{12pt}

\noindent
Retinopathy of Prematurity (ROP) is a vascular developmental disorder of the retina in
premature infants that can lead to blindness if not detected early. This study constructs a
classification pipeline for ROP staging from infant fundus images into five classes: Normal,
Stage 1, Stage 2, Stage 3, and laser scars. The study compares raw RGB images against a spatial
softmap representation designed to emphasize clinically relevant structures. The main
classification model is a custom residual CNN (TinyResNet) trained from scratch without
pretrained weights, receiving a three-channel input consisting of a Gabor-based vessel softmap,
a Meijering and oriented top-hat ridge softmap, and a CLAHE-normalized green channel. Group-aware
5-fold cross-validation yields a macro F1-score of \textbf{0.7853}. As a supporting evaluation,
the vessel segmentation workflow achieves a Dice score of \textbf{0.3881} on the HVDROPDB-BV
dataset, while the ridge channel reaches \textbf{0.0993} on HVDROPDB-RIDGE. A controlled ablation experiment
demonstrates that the softmap representation provides a +0.0987 improvement over raw RGB input
(0.6866 vs. 0.7853) under identical model, data, and protocol settings. These results confirm
that spatial vessel evidence---rather than raw pixel values---is the key driver of classification
performance improvement for ROP staging.

\vspace{12pt}

\noindent
Keywords: custom CNN, fundus image, Retinopathy of Prematurity, softmap representation,
TinyResNet
